%=============================================================================
% Wave 4, Iteration 18: Experimental Proposals --- Crisis Triage
% Agent 18: All experiments survive (lab-based)
%
% Model: Opus 4.6
% Date: 20 March 2026
%
% INHERITED STATE (from i17):
%   - SQMS Phase 0 READY ($47-50K, 2 weeks, blinding protocol)
%   - P7 Phase 0 READY ($1.55M, DOE+GA-EMS CRADA)
%   - TiN elevated to PRIMARY MEMS track
%   - Cat-psi 14-step protocol ($200-400K, 6 months)
%   - Modane combined demonstrator (4 patents, <$500K)
%   - mochi_class developer spec COMPLETE (19 functions, $40-80K)
%   - EXISTENTIAL: Dust-branch c_s ~ c/2.
%
% MANDATE (i18):
%   All experiments survive? Lab-based confirmation.
%=============================================================================

\documentclass[11pt]{article}
\usepackage{amsmath,amssymb}
\usepackage{geometry}
\geometry{margin=1in}
\usepackage{booktabs}
\usepackage{xcolor}
\usepackage{tcolorbox}

\title{\textbf{Wave 4, Iteration 18:}\\
\textbf{Experimental Proposals --- Crisis Triage}\\[6pt]
{\large All Lab Experiments Survive; Reprioritized for Sector A Confirmation}}
\author{DFD Research Programme --- W4i18 Agent Experimental (Opus 4.6)}
\date{20 March 2026}

\begin{document}
\maketitle

%=============================================================================
\section*{Executive Summary}
%=============================================================================

\begin{tcolorbox}[colback=green!5!white, colframe=green!70!black,
  title=\textbf{W4i18: ALL LAB EXPERIMENTS SURVIVE AND ARE REPRIORITIZED}]

Every lab-based experimental proposal is crisis-immune. The dust-branch
crisis INCREASES their importance because they become the primary
empirical path for DFD if cosmological predictions fail.

\begin{itemize}
\item \textbf{SURVIVES}: SQMS Phase 0, cat-psi, P7 Phase 0, MEMS/TiN,
  Modane demonstrator, LLR campaign. ALL crisis-immune.
\item \textbf{REPRIORITIZED}: \texttt{mochi\_class} elevated from
  ``important'' to ``existential.'' Must be funded immediately.
\item \textbf{AT RISK}: Euclid/DESI/SNe analysis protocols --- predictions
  may need revision pending \texttt{mochi\_class}.
\item \textbf{TOP FINDING}: Total lab programme cost $\sim$\$2.5M for
  comprehensive Sector A confirmation. This is unchanged by the crisis.
\end{itemize}
\end{tcolorbox}

%=============================================================================
\section{Experiment-by-Experiment Status}
%=============================================================================

\subsection{SQMS Phase 0 --- IMMUNE, HIGHEST PRIORITY}

\begin{itemize}
\item \textbf{What}: Q-ratio measurement in TESLA cavities.
\item \textbf{Cost}: \$47--50K. \textbf{Timeline}: 2 weeks.
\item \textbf{Discrimination}: $>30\sigma$ between DFD ($Q_{\rm ratio} =
  0.52$) and null ($Q_{\rm ratio} = 1.00$).
\item \textbf{Crisis impact}: NONE. Pure lab measurement.
\item \textbf{Crisis importance}: INCREASED. If cosmological sector fails,
  this is the first and fastest confirmation of Sector A.
\item \textbf{Status}: Ready for submission. Blinding protocol included.
\end{itemize}

\subsection{Cat-Psi Protocol --- IMMUNE}

\begin{itemize}
\item \textbf{What}: 14-step quantum protocol testing psi-modulation.
\item \textbf{Cost}: \$200--400K. \textbf{Timeline}: 6 months.
\item \textbf{Crisis impact}: NONE. Qubit decoherence is lab physics.
\item \textbf{Crisis importance}: INCREASED. Confirms psi-QEC (P3).
\end{itemize}

\subsection{P7 Energy Storage Phase 0 --- IMMUNE}

\begin{itemize}
\item \textbf{What}: Cavity energy storage demonstration.
\item \textbf{Cost}: \$1.55M (DOE+GA-EMS CRADA). \textbf{Timeline}: 1--2 yr.
\item \textbf{Crisis impact}: NONE. Cavity physics.
\item \textbf{Crisis importance}: INCREASED. If successful, provides
  commercial justification independent of all cosmology.
\end{itemize}

\subsection{MEMS/TiN Birefringence --- IMMUNE}

\begin{itemize}
\item \textbf{What}: TiN kinetic inductance readout for psi-field
  birefringence detection.
\item \textbf{Cost}: \$100--300K. \textbf{Timeline}: 6--12 months.
\item \textbf{Crisis impact}: NONE. Material property measurement.
\item \textbf{Crisis importance}: INCREASED. Simplest device-level test.
\end{itemize}

\subsection{Modane Combined Demonstrator --- IMMUNE}

\begin{itemize}
\item \textbf{What}: Underground facility testing 4 patents simultaneously.
\item \textbf{Cost}: $<$\$500K. \textbf{Timeline}: 1--2 years.
\item \textbf{Crisis impact}: NONE. All 4 tested patents are lab physics.
\item \textbf{Crisis importance}: UNCHANGED. Multi-patent efficiency.
\end{itemize}

\subsection{LLR at APOLLO --- IMMUNE}

\begin{itemize}
\item \textbf{What}: Lunar laser ranging for 0.93\,ns nonreciprocal.
\item \textbf{Cost}: Observing time only. \textbf{Timeline}: Months.
\item \textbf{Crisis impact}: NONE. Solar system $\psi$-field.
\item \textbf{Crisis importance}: INCREASED. Strongest single signal
  (SNR $\sim 4000$).
\end{itemize}

\subsection{mochi\_class --- EXISTENTIALLY REPRIORITIZED}

\begin{itemize}
\item \textbf{What}: Boltzmann code for DFD perturbation theory.
\item \textbf{Cost}: \$40--80K. \textbf{Timeline}: 2--3 months.
\item \textbf{Pre-crisis status}: Important for Euclid/DESI comparison.
\item \textbf{Post-crisis status}: \textbf{EXISTENTIAL}. Only tool that
  can determine whether DFD has a cosmological sector at all.
\item \textbf{Developer spec}: Complete (19 functions, 4 modules, 2
  CLASS hooks). Ready for implementation.
\end{itemize}

%=============================================================================
\section{Reprioritized Funding Sequence}
%=============================================================================

\begin{table}[h]
\centering
\begin{tabular}{clcc}
\toprule
\textbf{Priority} & \textbf{Experiment} & \textbf{Cost} & \textbf{Purpose} \\
\midrule
1 & SQMS Phase 0 & \$50K & Fastest Sector A test \\
2 & \texttt{mochi\_class} & \$60K & Existential crisis resolution \\
3 & Cat-psi at SQMS & \$300K & Deep Sector A confirmation \\
4 & MEMS/TiN & \$200K & Device-level confirmation \\
5 & Modane demonstrator & \$500K & Multi-patent efficiency \\
6 & P7 Phase 0 & \$1.55M & Commercial pathway \\
\midrule
& \textbf{TOTAL} & \textbf{\$2.66M} & \\
\bottomrule
\end{tabular}
\end{table}

\textbf{Key change}: \texttt{mochi\_class} elevated from position 5 to
position 2. It must run BEFORE Euclid DR1 (Oct 2026) to determine whether
the cosmological analysis protocols are valid.

%=============================================================================
\section{What Does NOT Change}
%=============================================================================

\begin{itemize}
\item Lab experiment designs: identical pre- and post-crisis.
\item Blinding protocols: unchanged.
\item Hardware requirements: unchanged.
\item Cost estimates: unchanged.
\item Personnel requirements: unchanged.
\end{itemize}

The crisis affects INTERPRETATION, not EXECUTION, of lab experiments.
All apparatus, protocols, and budgets proceed as designed.

%=============================================================================
\section*{TOP FINDING}
%=============================================================================

\begin{tcolorbox}[colback=blue!5!white, colframe=blue!70!black]
\textbf{All lab experiments are 100\% crisis-immune. Total Sector A
confirmation programme: \$2.66M. Only change: \texttt{mochi\_class}
elevated to priority \#2 (existential). The lab programme proceeds
unconditionally. If the cosmological sector fails, these experiments
become the ENTIRE empirical foundation of DFD --- making their execution
more important, not less.}
\end{tcolorbox}

\end{document}
